\section{Introduction}
\label{sec:intro}

High-frequency bus systems face a persistent coordination failure.
Timetables are planned offline to balance fleet utilization against passenger demand, yet the headways they prescribe degrade within minutes once buses encounter stochastic traffic and boarding delays.
Real-time holding control can restore regularity, but only within the headway structure it inherits.
When the timetable itself is poorly matched to demand---dispatching too few buses in the peak or too many in the off-peak---no amount of holding can recover acceptable service.
Planning and control are thus coupled: the quality of one layer bounds the achievable performance of the other.

Despite this coupling, existing work treats the two problems separately.
The timetable optimization literature formulates headway or frequency design as mixed-integer programs solved offline \citep{gkiotsalitis2021bus}, while the holding control literature trains reinforcement learning (RL) agents that take the timetable as given \citep{wang2020real}.
Joint optimization would seem a natural application of hierarchical RL, yet three gaps prevent a straightforward application.

\paragraph{Asynchrony between planning and control.}
Hierarchical RL frameworks such as the options framework \citep{sutton1999between} and goal-conditioned methods like HIRO \citep{nachum2018data} assume that the upper and lower levels share a common state clock or that the upper level sets subgoals consumed by the lower level within a fixed horizon.
In bus operations, the upper policy acts at dispatch events (once per bus departure), while the lower policy acts at every station arrival across all active buses.
The two levels observe different state spaces---the upper level sees route-level demand forecasts and fleet status, whereas the lower level sees per-bus headway deviations and passenger loads---and interact through delayed feedback spanning an entire trip lifecycle.
No shared clock or fixed subgoal horizon exists.

\paragraph{Cross-advantage under asynchrony.}
\citet{pan2024roboduet} recently introduced cross-advantage augmentation for bi-level robot locomotion, where an upper style policy and a lower tracking policy execute synchronously on the same state.
Their augmented advantage $\AUaug = \AU + \beta \AL$ directly sums the two advantages at each time step.
This elegant mechanism breaks down when the two levels operate on different timescales with different state and action spaces: there is no single time step at which $\AU$ and $\AL$ are jointly defined.
Generalising cross-level coupling to the asynchronous, multi-timescale setting requires propagating execution outcomes (rather than per-step advantages) over the variable-length trip lifecycle initiated by each dispatch, and exposing lower-level holding statistics to the upper policy through state augmentation rather than reward backflow.

\paragraph{Fleet and regularity constraints.}
Bus operations impose feasibility-style constraints---a fleet-size budget $\Nfleet$ and bounded headway deviation---that naive reward shaping handles poorly.
Constrained RL methods such as RCPO \citep{tessler2019reward} address single-level constraint satisfaction but do not account for constraints that span a bi-level structure, where the upper level controls resource allocation and the lower level controls service regularity.

\medskip

We propose TransitDuet, a bi-level RL framework that addresses all three gaps.
The upper policy $\piU$ outputs a per-dispatch target-headway shift $\delta_t \in [-120, 120]$\,s at each dispatch event; the lower's headway target for that trip is reset to $\htarget + \delta_t$. The lower policy $\piL$, a soft actor-critic with Lagrangian penalties \citep{haarnoja2018soft}, selects holding times at every station arrival with parameter sharing across buses, and its constraint is computed against the dispatch-specific goal headway.
This goal-conditioned coupling, in the spirit of HIRO~\citep{nachum2018data}, replaces the literal cross-advantage augmentation of RoboDuet~\citep{pan2024roboduet} (Section~\ref{sec:tap}); two auxiliary upper-only signals (a holding-feedback state channel and a per-dispatch hindsight credit on the upper reward) provide credit assignment without injecting upper signals into the lower's transitions.
Lower-policy training begins with a $30$-episode warmup during which the upper is held at $\delta_t = 0$, after which the goal channel engages.
Fleet-size feasibility is enforced by a $\thetavec$-weighted online gradient descent projection on the measurement vector $\zvec$ of dispatched buses, inspired by the ApproPO framework \citep{miryoosefi2019reinforcement}, while headway uniformity enters the lower-level objective through Lagrangian relaxation.

\paragraph{Contributions.}
The specific contributions of this paper are:
\begin{enumerate}
    \item \textbf{Goal-conditioned cross-level coupling.} We adopt a HIRO-style goal-conditioned coupling~\citep{nachum2018data}: the upper outputs a per-dispatch target-headway shift, and the lower's Lagrangian cost is computed against the resulting goal headway, with no upper-advantage injection into the lower's transitions. We compare this end-to-end against (i) a literal cross-advantage augmentation following HAAR~\citep{li2019hierarchical} with a PIPER-style reachability gate~\citep{singh2024piper}, and (ii) a three-channel coupling that uses $\delta_t$ as a launch-time perturbation with hindsight reward credit; the goal-conditioned design has the best mean composite cost ($1.444$ vs $1.499$ vs $1.492$) and the lowest seed variance ($1.7$--$2.0\times$ tighter than the alternatives) on the same evaluation protocol.
    \item \textbf{Joint timetable--holding formulation.} We cast bus timetable planning and holding control as a bi-level Markov decision process with heterogeneous state-action spaces and asynchronous transitions, providing a formal problem structure absent from prior bus RL work.
    \item \textbf{Constraint satisfaction via $\thetavec$-OGD and Lagrangian relaxation.} We integrate measurement-projection constraints at the upper level and dual-variable penalties at the lower level, handling fleet feasibility and headway regularity without reward shaping.
    \item \textbf{Target-policy-corrected lower SAC (TPC).} We identify a previously overlooked failure mode of joint bi-level RL on tightly coupled control: the lower policy is trained alongside an exploring upper policy, biasing the lower's training distribution away from deployment. We address this via importance-weighted lower SAC against an EMA "deployment" upper policy, with validation-based checkpoint selection.
    \item \textbf{Empirical evaluation on a calibrated corridor.} On a stochastic microsimulation under a unified 300-episode training and held-out evaluation protocol, TransitDuet attains the best mean on every metric we report: composite cost $1.444$ vs $1.504$ for the next-best baseline (CMA-ES, $-4.0\%$) and $1.544$ for Fixed ($-6.5\%$), with seed variance $1.7$--$2.5\times$ smaller than the search baselines.
\end{enumerate}

\paragraph{Results preview.}
We evaluate on a high-frequency bus corridor with 22 stations, 264 daily trips, demand stochasticity $\sigma_{\mathrm{demand}} = 0.15$, and elastic fleet $\Nfleet \in [8, 16]$ sampled per episode. Under a unified 300-episode training and held-out evaluation protocol, TransitDuet attains the best mean on wait, CV, overshoot, and composite cost simultaneously --- improvements of $5.5$/$6.4$/$10.4$/$6.5\%$ respectively over the strongest hand-tuned Fixed timetable --- and the lowest composite-cost variance of any method, while degrading only $25\%$ on wait when route stochasticity is doubled beyond the training distribution.

\paragraph{Paper organization.}
\Cref{sec:related} reviews related work in timetable optimization, bus holding RL, and hierarchical RL.
\Cref{sec:problem} formalizes the bi-level MDP.
\Cref{sec:method} presents the TransitDuet architecture, the cross-level coupling channels, and the constraint-handling mechanisms.
\Cref{sec:experiments} describes the simulation environment and reports results.
\Cref{sec:discussion} discusses limitations and extensions.
\Cref{sec:conclusion} concludes.
